\chapter*{Conclusion Générale}
\addcontentsline{toc}{chapter}{Conclusion Générale}

Au terme de ce projet, nous pouvons affirmer que l'objectif fixé au départ a été atteint~: doter
un centre de formation privé d'un outil décisionnel réel, fonctionnel et adapté à ses besoins
spécifiques.

Le constat de départ était simple mais révélateur~--- les responsables de centres de formation
naviguaient à vue, sans données consolidées ni indicateurs fiables pour prendre leurs décisions
au quotidien. Cette réalité nous a conduits à concevoir \textbf{MBICenter}, une plateforme qui
va bien au-delà d'un simple tableau de bord. Il s'agit d'un véritable système d'aide à la
décision, pensé pour accompagner chaque profil d'utilisateur dans ses responsabilités propres.

Le développement s'est articulé autour de trois \textit{releases} successives, conformément à
la méthodologie Scrum adoptée dès le début. La première a posé les fondations du système à
travers la gestion des accès et des utilisateurs. La deuxième a permis de structurer le c\oe{}ur
métier~--- sessions de formation, espace apprenant, alimentation des données. Enfin, la troisième
\textit{release} a constitué le véritable apport différenciant du projet~: des tableaux de bord
interactifs, des modèles prédictifs opérationnels et un module de recommandations stratégiques
piloté par intelligence artificielle générative.

Sur le plan technique, les choix effectués~--- Next.js, NestJS, FastAPI, Supabase~--- ont permis
de bâtir une architecture microservices robuste et évolutive. Les modèles de Machine Learning
intégrés ont affiché des performances satisfaisantes. Le recours à LLaMA~3 pour la
génération de recommandations personnalisées représente une avancée concrète vers une BI
augmentée par l'IA.

Ce projet nous a également permis de grandir en tant que développeuses~: gérer un cycle complet
de développement logiciel, collaborer en contexte professionnel réel au sein de Readdly Tech,
et surtout confronter nos choix techniques à des contraintes métiers concrètes.

Des pistes d'amélioration restent naturellement ouvertes~--- notamment l'enrichissement du module
de recommandations avec un mécanisme de validation et d'archivage, ou encore l'extension des
capacités prédictives à d'autres indicateurs. Mais dans son état actuel, MBICenter représente
une solution cohérente, utile et directement exploitable par les décideurs d'un centre de
formation.

En définitive, ce travail illustre comment les outils modernes de Business Intelligence,
combinés à l'intelligence artificielle, peuvent transformer des données brutes en leviers de
décision concrets~--- et contribuer à une gestion plus éclairée, plus réactive et tournée vers
l'avenir.